\section*{Chapter \ref{ch:s2:power-and-gas-trading} --- Power and Gas Trading}

\begin{solution}{exo:s2:power-and-gas-trading:1}
$1.5 \times 3.05 = 4.575$: about \euro4.6 per MWh below the day-ahead price, before other noise.
\end{solution}

\begin{solution}{exo:s2:power-and-gas-trading:2}
Charging 2~MWh of storage at \euro0 costs nothing; discharging returns $2\sqrt{0.88} = 1.876$~MWh to the grid (the charging side lost the other
$\sqrt{0.88}$ factor, paid at \euro0), sold at \euro100: \euro187.62.
\end{solution}

\begin{solution}{exo:s2:power-and-gas-trading:3}
$2.54 \times 0.61 \times 24 \times 727 / 2 \approx 13{,}500$: about \euro13,500 a year; with unrounded inputs, \euro13,600.
\end{solution}

\begin{solution}{exo:s2:power-and-gas-trading:4}
The edge is about \euro3.5 per MWh before costs (\euro2.54 after a \euro1 half-spread); a \euro3 half-spread takes most of it, leaving \euro0.54.
A forecast edge measured in a few euros is small against costs of a few euros.
\end{solution}

\begin{solution}{exo:s2:power-and-gas-trading:5}
Solar output peaks at midday; when it and wind exceed demand plus export capacity, producers that must or will run regardless (subsidised output,
inflexible plants) bid below zero. In 2024--2025, 23.4\% of hours at 13:00 were negative, none at 19:00.
\end{solution}

\begin{solution}{exo:s2:power-and-gas-trading:6}
They asked different questions with different data and models: Kiesel and Paraschiv let the response depend on a threshold in the expected demand
covered by conventional capacity, so the asymmetry can appear only in some states; Ziel tested whether positive and negative errors differ within his
model and found no significant difference. An effect that exists only above a threshold can vanish in an average.
\end{solution}

\begin{solution}{exo:s2:power-and-gas-trading:7}
Day-ahead revenue rises from \euro69,900 to \euro83,800, \euro13,800 more, and the intraday extra from \euro2,900 to \euro4,900. The table leaves out
degradation: each extra cycle uses up battery life, and the second cycle of the day uses the smaller price differences.
\end{solution}

\begin{solution}{exo:s2:power-and-gas-trading:8}
The backtest scheduled each day knowing that day's auction prices, which a real operator must forecast before bidding; it charged no degradation, no
bid-ask and no network charges. It is an upper bound for the day-ahead market alone, not a forecast; intraday and reserve revenues, which it omits,
push the other way.
\end{solution}

\begin{solution}{pb:s2:power-and-gas-trading:1}
\begin{enumerate}
\item The day-ahead price less the intraday price for the same hour, driven by what changed between the auction and delivery.
\item \euro3.05 per MWh per GW of residual load, correlation 0.83, over 2024--2025.
\item The day-ahead price less \euro3.05 per GW of unexpected wind, errors of 2~GW standard deviation with hourly autocorrelation 0.9, plus \euro6 of
noise; all of it planted.
\item Kiesel and Paraschiv: intraday prices adjust asymmetrically to renewable forecast errors and volume, depending on a demand threshold. Ziel: no
significant difference between the effects of positive and negative errors.
\item Sell day-ahead and buy back intraday when expecting more wind than the auction assumed; the reverse when less; nothing on small forecasts.
\item Skill 0.5, half-spread \euro1: \euro2.54 per MWh, \euro13,600 a year per MW, IC 0.36; skill 0.3: \euro1.08, \euro5,800, 0.21; half-spread
\euro3: \euro0.54, \euro2,900.
\item A position in power or gas on a weather forecast that differs from the market's, closed when forecasts converge or the weather arrives.
\item Better public forecasts and costs.
\item Scheduling a battery to charge cheap and discharge dear, day-ahead then intraday, within its limits.
\item Cheap at midday, dear in the evening; 457 negative hours in 2024 and 573 in 2025, 23.4\% of hours at 13:00.
\item Daily dynamic programming over state of charge and energy charged, one cycle, 88\% round trip, ending empty.
\item Day-ahead \euro69,900 a year; intraday extra \euro2,900 (\euro3,800 without spread, \euro1,600 at \euro3); two cycles \euro83,800 and
\euro4,900.
\item \euro69,900 per MW a year day-ahead, \euro72,800 with intraday re-optimisation (\euro2,900 more), on 2024--2025 German prices.
\item The day-ahead schedule knows the auction's prices; the intraday one knows the day's intraday prices at once.
\item The planted errors move prices by a few euros, while the daily range averages over a hundred.
\item Correlated forecast errors across hours, price jumps, imbalance charges, the asset's limits and degradation.
\item Bid on price forecasts made before the auction, re-optimise intraday as prices arrive, and charge degradation and fees.
\item Battery day-ahead arbitrage: more batteries flatten the daily shape it trades.
\item Both are paid by participants who must trade (flows at a fix, generators balancing errors), and both need a forecast of that need.
\item Delivery at a particular hour, and the information about what that hour will need.
\end{enumerate}
\end{solution}

\begin{solution}{iq:s2:power-and-gas-trading:1}
Some producers are paid to produce (subsidies per MWh) or cannot stop cheaply (inflexible plants), so they bid below zero rather than cut output; when
their supply exceeds demand and export capacity, the auction clears below zero.
\end{solution}

\begin{solution}{iq:s2:power-and-gas-trading:2}
Forecast the change in wind, solar and load between the auction's forecasts and one's own, translate it to price through the residual-load slope in
that hour's state of the merit order, and check against realised spreads; include plant outages and interconnector flows.
\end{solution}

\begin{solution}{iq:s2:power-and-gas-trading:3}
Sell those hours day-ahead and plan to buy back intraday, sized by the forecast's track record at that lead time and the expected price effect
(about \euro3 per GW here), and within limits for correlated errors across the afternoon's hours.
\end{solution}

\begin{solution}{iq:s2:power-and-gas-trading:4}
Limits on open position by hour and by block, on the sum of correlated hours, on imbalance exposure at gate closure, on loss per day, and stress tests
for price spikes and forecast failures.
\end{solution}

\begin{solution}{iq:s2:power-and-gas-trading:5}
A price-forecast service; an optimiser per battery (dynamic programming or a linear programme) re-run on each new price; a position keeper that
nets the fleet against markets; order routing to the intraday market; telemetry of state of charge and availability; and a fallback schedule.
\end{solution}

\begin{solution}{iq:s2:power-and-gas-trading:6}
Buying 1~MWh at $p_1$ returns $\eta$~MWh sold at $p_2$: profit $\eta p_2 - p_1 > 0$ iff $p_2 > p_1/\eta$. With $p_1 < 0$ the battery is paid to
charge, so the trade pays whenever $\eta p_2 > p_1$, including some negative $p_2$; the constraint then is capacity, not price.
\end{solution}
